% problem_id: S046-thm:1surj_to_surj_to_pem
% source_id: S046 (HoTT Book)
% source_file: setmath.tex
% statement_locator: setmath.tex lines 614--619
% proof_locator: setmath.tex lines 621--646
% env: thm
% pairing: tight (verified)
% status: complete (statement + author proof verbatim + reason + locators)
%
\begin{thm}[Diaconescu]\label{thm:1surj_to_surj_to_pem}
  \index{axiom!of choice}%
  \index{excluded middle}%
  \index{Diaconescu's theorem}\index{theorem!Diaconescu's}%
  The axiom of choice implies the law of excluded middle.
\end{thm}

\begin{proof}
We use the equivalent form of choice given in \cref{thm:ac-epis-split}.
Consider a mere proposition $A$.
The function $f:\bool\to\susp(A)$ defined by
$f(\bfalse) \defeq \north$ and $f(\btrue) \defeq \south$
is surjective.
Indeed, we have
$\pairr{\bfalse,\refl{\north}} : \hfiber{f}{\north}$
and $\pairr{\btrue,\refl{\south}} :\hfiber{f}{\south}$.
Since $\bbrck{\hfiber{f}{x}}$ is a mere proposition, by induction the claimed surjectivity follows.

By \cref{prop:trunc_of_prop_is_set} the suspension $\susp(A)$
is a set, so by the axiom of choice there merely exists a
section $g: \susp(A) \to \bool$ of $f$.
As equality on $\bool$ is decidable we get
\begin{equation*}
 (g(f(\bfalse))= g(f(\btrue))) +
 \lnot (g(f(\bfalse))= g(f(\btrue))),
\end{equation*}
and, since $g$ is a section of $f$, hence injective,
\begin{equation*}
(f(\bfalse) = f(\btrue)) +
\lnot (f(\bfalse) = f(\btrue)).
\end{equation*}
Finally, since $(f(\bfalse)=f(\btrue)) = (\north=\south) = A$ by \cref{prop:trunc_of_prop_is_set}, we have $A+\neg A$.
\end{proof}
